% Chapter 3 -- The Detection Pipeline
% Standalone, compilable draft (FPT-template article class + booktabs + tikz).
% Folded into the thesis main document later via \input.
%
% Honesty-phrasing rules applied (personal_agenda P1):
%   - production classifier is the FROM-SCRATCH Stage B checkpoint
%     (stage_b_scratch_best.pth, Finding #31); synthetic Stage A pretraining
%     was DROPPED from the final pipeline and is kept only as Ch-4 ablation
%     material (#25/#30). We do NOT present two-stage training as a contribution.
%   - measured vs inferred evidence tiers marked in prose; timing numbers are
%     measured (timing_seq08_full_n4068.json), cited as such.
%   - every number traceable to a finding / output artifact.
%   - "the classifier's job is precision" claim deferred to the Ch-4 ablation.
%
\chapter{Methodology}
\label{ch:pipeline}

This chapter describes the method in the order a scan flows through it: a modular,
classical car-detection pipeline (Section~\ref{sec:meth-detection}) that segments
and classifies car clusters, an object-completion stage
(Section~\ref{ch:completion-method}) that reconstructs each detected car's full
surface from the one-sided partial a single frame observes, and the reference-free
donor metric (Section~\ref{sec:donor-metric}) that makes that completion measurable
on real data. The completion is evaluated with this metric in
Chapter~\ref{ch:completion-eval}.

\section{Detection pipeline}
\label{sec:meth-detection}

This section describes the car-detection pipeline that produces every detection
evaluated in Chapter~\ref{ch:detection-eval} and every cluster completed in
Section~\ref{ch:completion-method} and Chapter~\ref{ch:completion-eval}. The
design is deliberately modular and classical: a chain of independent stages,
each with a small number of interpretable parameters, rather than an end-to-end
learned detector. This choice trades peak accuracy for interpretability,
reproducibility, and near-zero manual annotation --- the classifier is trained on
labels mined automatically from the pipeline itself
(Section~\ref{sec:pipeline-classifier}) --- and it is what lets Chapter~\ref{ch:recall-bottleneck}
root-cause the recall shortfall to a single stage rather than to an opaque
network.

The pipeline takes one raw LiDAR scan and emits a set of car detections, each an
oriented bounding box with a stable track identity across frames. Every stage
parameter lives in a single dictionary, \texttt{PIPELINE\_CONFIG} in
\texttt{src/pipeline.py}, which was frozen before the write-up began; the frozen
values are collected in Table~\ref{tab:pipeline-config}. Figure~\ref{fig:pipeline}
gives the full chain with the measured per-stage cost.

\subsection{Overview}
\label{sec:pipeline-overview}

\begin{figure}[H]
  \centering
  \begin{tikzpicture}[
      font=\small,
      node distance=6mm,
      stage/.style={
        draw, rounded corners, align=center,
        minimum width=32mm, minimum height=11mm,
        fill=blue!5, thick},
      learned/.style={stage, fill=orange!12},
      io/.style={draw, ellipse, align=center, fill=black!5,
        minimum width=26mm, minimum height=9mm},
      ms/.style={font=\footnotesize\itshape, text=black!60},
      arr/.style={-{Stealth[length=2.2mm]}, thick},
    ]
    \node[io] (raw) {Raw LiDAR scan\\$\sim$120k points};
    \node[stage, below=of raw] (pre) {Preprocess\\{\footnotesize $z$-filter $\to$ voxel $\to$ denoise}};
    \node[stage, below=of pre] (ground) {Ground removal\\{\footnotesize RANSAC plane}};
    \node[stage, below=of ground] (cluster) {Clustering\\{\footnotesize HDBSCAN @ \SI{0.10}{\meter}}};
    \node[stage, below=of cluster] (geo) {Geometric filter\\{\footnotesize OBB gates}};
    \node[learned, below=of geo] (clf) {Classifier\\{\footnotesize PointNet, car/not-car}};
    \node[stage, below=of clf] (track) {Tracker\\{\footnotesize centroid + track vote}};
    \node[io, below=of track] (out) {Car detections\\{\footnotesize boxes + track ids}};
    \node[learned, right=18mm of track] (comp) {Completion\\{\footnotesize PCN (Ch.~6)}};

    \draw[arr] (raw) -- (pre);
    \draw[arr] (pre) -- (ground);
    \draw[arr] (ground) -- (cluster);
    \draw[arr] (cluster) -- (geo);
    \draw[arr] (geo) -- (clf);
    \draw[arr] (clf) -- (track);
    \draw[arr] (track) -- (out);
    \draw[arr] (track) -- (comp);

    % per-stage timing annotations (right of each stage)
    \node[ms, right=2mm of pre]     {\SI{126}{\milli\second} (20\%)};
    \node[ms, right=2mm of ground]  {\SI{56}{\milli\second} (9\%)};
    \node[ms, right=2mm of cluster] {\SI{370}{\milli\second} (59\%)};
    \node[ms, right=2mm of geo]     {\SI{36}{\milli\second} (6\%)};
    \node[ms, below=0.5mm of comp.south, text width=34mm, align=center]
      {\footnotesize\itshape +\SI{13}{\milli\second} per completed car};
    \node[ms, left=2mm of clf]      {\SI{34}{\milli\second} (5\%)};
    \node[ms, left=2mm of track]    {\SI{0.3}{\milli\second}};
  \end{tikzpicture}
  \caption{The detection pipeline. Learned stages are shaded; all others are
    classical. Per-stage figures are mean wall-clock time on the full
    sequence~08 under the final (promoted) configuration
    (\num{4068} frames, AMD Ryzen 7 7800X3D + RTX 3070~Ti;
    \texttt{output/experiments/timing/timing\_seq08\_full\_n4068\_combined.json}), and are
    measured, not estimated. The total is
    \SI{623.5}{\milli\second\per\frame} ($\approx$\SI{1.6}{\fps}); the only learned
    \emph{per-frame} stage --- the classifier --- accounts for just
    \SI{34}{\milli\second} (\SI{5}{\percent}), and the classical clustering and
    ground-removal stages dominate. Completion runs off the tracked detections and
    is optional; it is learned, but its cost is per completed car, not per frame, so
    it does not enter this per-frame budget.}
  \label{fig:pipeline}
\end{figure}

The system is \emph{offline} by design: the tracker and the track-level filter
use the whole sequence, and real-time operation is out of scope. At
\SI{623.5}{\milli\second\per\frame} the pipeline is roughly \SI{1.6}{\fps}, an order
of magnitude below sensor rate; we characterise this cost honestly in the runtime
section of Chapter~\ref{ch:detection-eval} rather than claiming real-time capability. The one
observation worth making here is where the cost lives: clustering alone is
\SI{59}{\percent} of the frame budget, and the only learned per-frame stage, the
classifier, is just \SI{5}{\percent}. The expensive parts of this pipeline are
classical, not neural.

\subsection{Preprocessing}
\label{sec:pipeline-preprocess}

Each raw scan (roughly \num{120000} points) passes through three cleaning steps
before any object reasoning:

\begin{enumerate}
  \item \textbf{Ground-height $z$-filter.} Points below
    $z = \SI{-2.0}{\meter}$ in the sensor frame are dropped. This removes the bulk
    of the road surface and any sub-ground noise cheaply, before the more
    expensive steps see them.
  \item \textbf{Voxel downsampling} at \SI{0.05}{\meter}. Points are collapsed to
    one representative per \SI{5}{\centi\meter} voxel, which both regularises the
    range-dependent density of LiDAR and cuts the point count by roughly an order
    of magnitude.
  \item \textbf{Statistical outlier removal.} Points whose mean distance to their
    \num{20} nearest neighbours exceeds \num{2.0} standard deviations of the
    local distribution are discarded, removing isolated speckle.
\end{enumerate}

The order of the last two steps is behaviour-changing. In the frozen configuration
(\texttt{voxel\_\allowbreak before\_\allowbreak denoise = True}) downsampling runs \emph{before}
denoising, so the statistical outlier test operates on roughly ten times fewer
points --- and on a regularised point set, which changes which points survive.
This ordering was adopted as part of the Finding~\#34 optimisation primarily
because it improves detection (discussed in Chapter~\ref{ch:detection-eval}); its effect on the measured
denoising cost is small. It is behaviour-changing, so we flag it here rather than
presenting the two orderings as equivalent.
Preprocessing costs \SI{126}{\milli\second\per\frame} in total, split across the
$z$-filter (\SI{58}{\milli\second}), denoising (\SI{50}{\milli\second}), and
downsampling (\SI{18}{\milli\second}).

\subsection{Ground removal}
\label{sec:pipeline-ground}

The dominant remaining structure after preprocessing is the ground plane, which
must be removed before clustering or it bridges every object into one component.
We fit a single plane with RANSAC: over \num{300} iterations, three points are
sampled, their plane computed, and inliers counted within
\SI{0.2}{\meter}; the best-supported plane is kept. The plane is
\emph{accepted only if its normal is near-vertical} --- the $z$-component of the
unit normal must exceed \num{0.5} --- so a spuriously-fitted vertical surface (a
wall, the side of a truck) cannot be mistaken for ground. If the check fails, no
points are removed and the whole cloud is passed on as objects; this fail-safe is
rare but prevents a bad plane fit from deleting real cars.

A single global plane is a deliberate simplification: it assumes the local ground
is approximately planar, which holds on the flat urban sequences used here but
would degrade on strong slopes. A grid-based ground model (Patchwork++-style) is
listed as future work, not a limitation we claim to have solved. The iteration
count was reduced from \num{1000} to \num{300} as part of Finding~\#34 after
confirming the plane fit converges well before \num{1000} iterations on this
data; the fit uses a fixed seed (\texttt{np.random.default\_rng(42)}) so ground
removal, and therefore the whole pipeline, is deterministic across runs. Ground
removal costs \SI{56}{\milli\second\per\frame} (\SI{9}{\percent} of the frame).

\subsection{Clustering}
\label{sec:pipeline-cluster}

Non-ground points are grouped into object candidates with HDBSCAN, a
density-based hierarchical clusterer, using a minimum cluster size of \num{10}
and minimum samples of \num{5}. HDBSCAN is chosen over a fixed-radius method
because LiDAR density falls off with range: a single Euclidean radius that merges
nothing at \SI{5}{\meter} will over-segment at \SI{40}{\meter}. HDBSCAN's
variable-density hierarchy adapts to this without a hand-tuned radius schedule,
and Chapter~\ref{ch:detection-eval} reports a controlled benchmark (Finding~\#43) confirming it beats
DBSCAN and Euclidean extraction on recall at matched settings.

Two details of the frozen configuration are load-bearing. First, clustering runs
on a \emph{coarser} \SI{0.10}{\meter} voxel grid than the \SI{0.05}{\meter}
preprocessing grid: the object cloud is re-downsampled to \SI{0.10}{\meter},
HDBSCAN is run there, and the labels are propagated back to the full-resolution
points by nearest neighbour. This was adopted (Finding~\#34) because the coarser
grid closes intra-car density gaps that otherwise cause a single car to split
into several clusters --- the central mechanism of the recall bottleneck analysed
in Chapter~\ref{ch:recall-bottleneck} --- and it added \num{1655} true positives on the full sequence~08.
Second, HDBSCAN is by far the most expensive stage at
\SI{370}{\milli\second\per\frame} (\SI{59}{\percent} of the budget), because it is
superlinear in the point count; the coarse-grid trick is as much a speed measure
as an accuracy one. The clustering stage produces on average \num{45} candidate
clusters per frame.

\subsection{Geometric filtering}
\label{sec:pipeline-geometric}

Most clusters are not cars: they are trees, poles, building facades, and
fragments. A geometric filter rejects clusters whose oriented bounding box (OBB)
is inconsistent with a car, using only cheap, interpretable shape gates. For each
cluster the OBB is computed and the cluster is kept only if it passes \emph{all}
of the following:

\begin{itemize}
  \item at least \num{15} points;
  \item OBB volume in $[\num{0.5}, \num{50}]\,\si{\meter\cubed}$;
  \item longest OBB extent $\leq \SI{6.0}{\meter}$, and the two largest extents
    above small floors (\SI{0.5}{} and \SI{0.4}{\meter}), rejecting slivers;
  \item box centre no more than \SI{1.5}{\meter} above the fitted ground plane,
    and vertical point span $\leq \SI{1.8}{\meter}$, rejecting tall structures;
  \item aspect ratio (longest\,/\,shortest extent) $\leq \num{6.0}$, rejecting
    thin walls and poles.
\end{itemize}

These bounds are generous car priors, not a tight car model --- they are meant to
remove obvious non-cars while leaving the fine car/not-car decision to the
learned classifier. The height gate is computed \emph{relative to the fitted
ground plane}, not to a fixed world height, so it is robust to the sensor not
being perfectly level. The filter costs \SI{36}{\milli\second\per\frame}
(\SI{6}{\percent}). Its individual gates are ablated in Chapter~\ref{ch:detection-eval}; here we note
only that this stage is high-recall and low-precision by construction, and the
classifier is what turns its survivors into precise detections.

\subsection{Classifier}
\label{sec:pipeline-classifier}

Each surviving cluster is classified as \emph{car} or \emph{not-car} by a small
dual-branch PointNet (\texttt{src/\allowbreak classifier.py}). The binary decision
replaces
an earlier multi-class (car/bus/motorcycle) design that was dropped once the data
showed near-zero real bus and motorcycle instances (Finding~\#20); the scope
narrowing to cars is deliberate and is owned as such in the introduction.

\paragraph{Architecture.} The network has two branches
(Figure~\ref{fig:classifier}) and about \num{80000}
parameters --- deliberately small (Finding~\#12):
\begin{itemize}
  \item a \textbf{point branch} that takes \num{512} points (sampled or padded)
    normalised into the unit sphere, and applies shared per-point MLPs
    ($3 \to 64 \to 128 \to 256$) followed by a max-pool to a \num{256}-d shape
    descriptor;
  \item a \textbf{bounding-box branch} that takes an \num{8}-dimensional metric
    feature vector (sorted OBB extents, volume, two aspect ratios,
    $\log(1+\text{point count})$, vertical span), $z$-score normalised, through a
    small MLP to \num{32}-d.
\end{itemize}
The two are concatenated and passed through a small head
($288 \to 128 \to 2$) with dropout. The bounding-box branch is what re-injects
absolute metric scale, which the unit-sphere normalisation deliberately strips
from the point branch: a distant small car and a nearby one have the same
normalised shape but different real extents, and the metric branch is how the
network sees the difference. There are no T-Nets, since clusters already sit in a
consistent sensor frame.

\begin{figure}[H]
  \centering
  \begin{tikzpicture}[
      font=\footnotesize, node distance=6mm and 7mm,
      blk/.style={draw, rounded corners, align=center, fill=blue!5, thick,
        minimum height=8mm, inner sep=3pt},
      lrn/.style={blk, fill=orange!12},
      io/.style={draw, ellipse, align=center, fill=black!5, inner sep=2pt},
      arr/.style={-{Stealth[length=2mm]}, thick},
    ]
    \node[io] (cl) at (0,0) {cluster};
    \node[blk] (pts) at (2.3,1.3) {\num{512} points\\unit-sphere};
    \node[lrn, right=of pts] (pmlp) {shared MLP\\{\scriptsize $3\!\to\!64\!\to\!128\!\to\!256$}};
    \node[blk, right=of pmlp] (pool) {max-\\pool};
    \node[blk, right=of pool] (pd) {shape\\\num{256}-d};
    \node[blk] (feat) at (2.3,-1.3) {\num{8} OBB\\features};
    \node[lrn, right=of feat] (bmlp) {MLP};
    \node[blk, right=of bmlp] (bd) {metric\\\num{32}-d};
    \node[blk] (cat) at ($(pd.east)+(1.1,-1.3)$) {concat\\\num{288}-d};
    \node[lrn, right=of cat] (head) {head\\{\scriptsize $288\!\to\!128\!\to\!2$}};
    \node[io, right=of head] (out) {car /\\not-car};
    \draw[arr] (cl) |- (pts.west);
    \draw[arr] (cl) |- (feat.west);
    \draw[arr] (pts)--(pmlp); \draw[arr] (pmlp)--(pool); \draw[arr] (pool)--(pd);
    \draw[arr] (feat)--(bmlp); \draw[arr] (bmlp)--(bd);
    \draw[arr] (pd) -| (cat.north);
    \draw[arr] (bd) -| (cat.south);
    \draw[arr] (cat) -- (head);
    \draw[arr] (head) -- (out);
  \end{tikzpicture}
  \caption{The dual-branch classifier (\texttt{src/classifier.py}), applied to one
    cluster. The \emph{point branch} takes \num{512} unit-sphere-normalised points
    through shared per-point MLPs and a max-pool to a \num{256}-d shape descriptor;
    the \emph{bounding-box branch} takes \num{8} metric shape features --- the
    absolute scale the sphere normalisation strips --- to \num{32}-d. The two are
    concatenated and passed through a small head to a car/not-car decision. Learned
    blocks are shaded.}
  \label{fig:classifier}
\end{figure}

\paragraph{Training data (mined, not hand-labelled).} The classifier is trained
on real LiDAR clusters mined automatically from SemanticKITTI
(\texttt{src/mine\_stage\_b.py}). For each training frame the pipeline's own
stages~1--5 are replayed, the ground-truth semantic labels are propagated onto the
surviving clusters by nearest neighbour, and each cluster is labelled \emph{car}
or \emph{not-car} by majority vote, keeping only clusters whose majority class
has purity $\geq \num{0.75}$. This yields \num{420333} training clusters
(sequences 00--07, 09, 10) and \num{130394} validation clusters (sequence~08),
with the intrinsic $\approx$21\,\%/79\,\% car/not-car imbalance handled by a
class-weighted loss. The key property is that the training clusters are produced
by the \emph{same} clustering and filtering the classifier meets at inference ---
including HDBSCAN splits and filter survivors --- so it learns on exactly the
distribution it must classify, at no manual-annotation cost.

\paragraph{Production checkpoint: trained from scratch.} The production classifier
is \texttt{stage\_b\_\allowbreak scratch\_\allowbreak best.pth}, trained on the
mined real clusters from
random initialisation (val macro-$F_1$ \num{0.9285}, Finding~\#31). An earlier
two-stage design first pretrained on synthetic ShapeNet partial renders
(``Stage~A'') and then fine-tuned on the real clusters (``Stage~B''). We report
the two-stage machinery in Chapter~\ref{ch:detection-eval} as an \emph{ablation}, because a controlled
comparison (Finding~\#25) showed the synthetic pretraining to be \emph{redundant}
given \num{420000} real clusters: the from-scratch classifier matches the
pretrained one (macro-$F_1$ \num{0.9285} vs.\ \num{0.9225}). The final pipeline
therefore uses the simpler from-scratch checkpoint, and the synthetic
pretraining, together with the finding that the sim-to-real gap is total and
symmetric (Finding~\#30), is presented as a negative result rather than as a
contribution. The classifier costs \SI{34}{\milli\second\per\frame} (about
\SI{0.75}{\milli\second} per cluster over the \num{45} clusters in a frame).

\subsection{Tracking and track-level filtering}
\label{sec:pipeline-tracker}

Detections are linked across frames by a centroid tracker
(\texttt{src/tracker.py}). Each frame, current detections are matched to existing
tracks by solving a Hungarian assignment on the pairwise centroid-distance
matrix; matches beyond \SI{2.0}{\meter} are rejected, unmatched detections start
new tracks, and a track that goes unmatched for more than \num{5} consecutive
frames is retired. This is an offline, centroid-only tracker --- no motion model,
no appearance --- which is adequate here because detection, not association, is
the bottleneck; an IoU-based (SORT-style) tracker is listed as future work.

The tracker also enables a \emph{track-level filter} that improves precision.
Each detection carries the classifier's per-frame label; a track is accepted as a
car only if, across its lifetime, it collected at least \num{2} non-rejected
votes, those known votes are a majority of its votes
($\geq \num{0.5}$), and they agree on a single class. A cluster that the
classifier flickers on --- a car in one frame, not-car in the next --- is thus
resolved by temporal majority rather than trusted per frame. Because this filter
consults the whole track, it is an offline post-processing step, consistent with
the offline scope of the system; its contribution is isolated as an ablation row
in Chapter~\ref{ch:detection-eval}. The tracker itself is effectively free at
\SI{0.3}{\milli\second\per\frame}.

\subsection{Frozen configuration}
\label{sec:pipeline-config}

Table~\ref{tab:pipeline-config} lists the frozen parameter values that produce
every number in this thesis. They are the committed \texttt{PIPELINE\_CONFIG} and
were not altered to obtain any reported result.

\begin{table}[H]
  \centering
  \caption{Frozen pipeline configuration (\texttt{src/pipeline.py}). These values
    were fixed before the write-up; no reported number was produced by changing
    them.}
  \label{tab:pipeline-config}
  \begin{tabular}{lll}
    \toprule
    Stage & Parameter & Value \\
    \midrule
    Preprocess & $z$-filter threshold & \SI{-2.0}{\meter} \\
               & voxel size & \SI{0.05}{\meter} \\
               & denoise neighbours / std-ratio & \num{20} / \num{2.0} \\
               & voxel before denoise & true \\
    \midrule
    Ground     & RANSAC iterations & \num{300} \\
               & inlier distance & \SI{0.2}{\meter} \\
               & min.\ normal $z$ & \num{0.5} \\
    \midrule
    Clustering & method & HDBSCAN \\
               & min.\ cluster size / min.\ samples & \num{10} / \num{5} \\
               & clustering voxel size & \SI{0.10}{\meter} \\
    \midrule
    Geometric  & min.\ points & \num{15} \\
               & volume range & $[\num{0.5}, \num{50}]\,\si{\meter\cubed}$ \\
               & max.\ longest extent & \SI{6.0}{\meter} \\
               & max.\ height above ground & \SI{1.5}{\meter} \\
               & max.\ vertical span & \SI{1.8}{\meter} \\
               & max.\ aspect ratio & \num{6.0} \\
    \midrule
    Classifier & checkpoint & \texttt{stage\_b\_scratch\_best.pth} \\
               & input points & \num{512} \\
    \midrule
    Tracker    & max.\ match distance & \SI{2.0}{\meter} \\
               & max.\ disappeared frames & \num{5} \\
               & min.\ track length & \num{2} \\
               & min.\ known votes / ratio & \num{2} / \num{0.5} \\
    \bottomrule
  \end{tabular}
\end{table}

The completion stage that operates on the tracked car detections
(Figure~\ref{fig:pipeline}, right) is the subject of Section~\ref{ch:completion-method}; the evaluation of
the detections themselves is Chapter~\ref{ch:detection-eval}.
